%-------------------------------------------------------------------------------
%	SECTION TITLE
%-------------------------------------------------------------------------------
\cvsection{İş~Deneyimi}
\vspace{-2mm}

%-------------------------------------------------------------------------------
%	CONTENT
%-------------------------------------------------------------------------------
\cvsubsection{İNCİ GS YUASA/İNCİ AKÜ}
\vspace{-2mm}

\begin{cventries}

%---------------------------------------------------------
  \cventry
    {OPERASYONEL TEKNOLOJİLER EKİP LİDERİ --- 3 Üretim Tesisi, Türkiye} % Job title
    {} % Organization
    {} % Location
    {Oca 2026 - Devam Ediyor} % Date(s)
    {
      \begin{cvitems} % Description(s) of tasks/responsibilities
        \item {\textbf{3 üretim tesisinde} OT dijitalleşmesi ve uygulamalı yapay zeka stratejisine liderlik ediyor; 2 kişilik bir mühendislik ekibini yönetiyor ve \textbf{10'dan fazla} fonksiyon/departmanla koordinasyon sağlıyor. Göreve başladığından bu yana \textbf{15+} süreç iyileştirme projesi teslim etti.}
        \item {OT veri entegrasyon altyapısını \textbf{PLC}, \textbf{SCADA}, \textbf{MES} ve historian sistemleri genelinde \textbf{300+} yönetilen veri kanalı ve veritabanı bağlantısına ölçeklendirdi (yaklaşık 200'den yükseldi) --- 2026 yılında tek başına \textbf{36} yeni üretim hattı/makine (fırınlar, değirmenler, montaj, dolum ve kaynak istasyonları) entegre etti ve herhangi bir kullanıcının tüm bağlı makineler için gerçek zamanlı özel alarmlar yapılandırabildiği kendi kendine hizmet veren \textbf{OPC UA} tabanlı bir uyarı platformu kurdu.}
        \item {Uygulamalı yapay zeka portföyünü \textbf{4} aktif kestirimci bakım/kalite modeline ve \textbf{25} yapay zeka destekli otomasyon projesine büyüttü; günlerce süren raporlama döngülerini saniyelere indiren RAG tabanlı bilgi sistemleri ile GenAI/çoklu ajan iş akışları (\textbf{n8n}, Microsoft Copilot, GPT, Claude) dahil.}
        \item {Yapay zeka destekli insan ve forklift algılama güvenlik sistemi, yüksek maruziyetli alanlarda gerçek zamanlı hava kalitesi/partikül izlemesi (çalışanların kandaki kurşun seviyelerinde ölçülebilir azalmaya katkı sağladı) ve potansiyel bir yangın olayını önlemeye yardımcı olan üretim/yardımcı tesis sistemlerinde gerçek zamanlı anomali tespiti devreye aldı --- genişleyen bir uyum programının parçası olarak (denetim aksiyonları yıllık 3'ten 7'ye yükseldi) \textbf{4 yıl kesintisiz sıfır iş kazası} sürdürdü.}
        \item {OT fonksiyonu için Ar-Ge ve inovasyon görünürlüğüne liderlik ediyor --- \textbf{TÜBİTAK} destekli bir Ar-Ge projesini koordine ediyor ve yıllık \textbf{6} iç/dış inovasyon yarışması başvurusu hazırlıyor; birden fazla birincilik ve İnovasyon Özel Ödülü kazanımına katkı sağladı.}
      \end{cvitems}
    }

%---------------------------------------------------------
  \cventry
    {KIDEMLİ OPERASYONEL TEKNOLOJİLER (OT) MÜHENDİSİ --- Manisa, Türkiye} % Job title
    {} % Organization
    {} % Location
    {Ara 2022 - Ara 2025} % Date(s)
    {
      \begin{cvitems} % Description(s) of tasks/responsibilities
        \item {\textbf{ISA95} standartlarına dayalı bir \textbf{Unified Namespace (UNS)} uygulaması ile üretim çevrim sürelerini azaltarak operasyonel verimlilikte, iyileşen üretim verimiyle ölçülen \textbf{\%19} artış sağladı.}
        \item {\textbf{OPC UA}, \textbf{MQTT} ve \textbf{Kafka} ile altyapı optimizasyonu yoluyla gerçek zamanlı veri akışını geliştirerek veri güvenilirliğinde \textbf{\%14} iyileşme sağladı; bu da müşteri şikayetlerine yanıt süresinde \%90 iyileşmeye yol açtı.}
        \item {Kurşun-asit akü üretimi için \textbf{Grafana} ve \textbf{Node.js} ile gerçek zamanlı panolar tasarlayarak üretim izleme gecikmelerini, üretim anomalilerinin daha hızlı tespit edilmesiyle ölçülen \textbf{\%32} oranında azalttı.}
        \item {\textbf{OT/ICS} sistemlerini modernize edip otomasyon teknolojilerini entegre ederek maliyet tasarrufu sağlayan dijital dönüşüm girişimlerine liderlik etti; azalan genel giderlerle yansıyan operasyonel giderlerde \textbf{\%57} azalma sağladı.}
      \end{cvitems}
    }

%---------------------------------------------------------
\end{cventries}

\cvsubsection{DVD Valves --- Manisa, Türkiye}
\vspace{-2mm}

\begin{cventries}

%---------------------------------------------------------
  \cventry
    {YAZILIM GELİŞTİRME LİDERİ} % Job title
    {} % Organization
    {} % Location
    {Oca 2022 - Kas 2022} % Date(s)
    {
      \begin{cvitems} % Description(s) of tasks/responsibilities
        \item {Akıllı sulama hidrantları için gömülü sistemler ve \textbf{web servisi} entegrasyonunu yöneten bir TEYDEB projesine (\$200k finansman) \textbf{Node.js} ve \textbf{React} kullanarak liderlik etti; raporlama ve muhasebe hizmetleri verimliliğinde \textbf{\%22} iyileşme sağladı.}
        \item {\textbf{DCS} kullanarak sulama için kontrol sistemleri geliştirdi; optimize edilmiş sistem yanıt süreleriyle kanıtlanan \textbf{\%18} verimlilik artışı sağladı.}
        \item {Büyük ölçekli uluslararası su projelerinde (\textbf{İtalya}, \textbf{Dubai}, \textbf{Kolombiya}) hidrolik kelebek vanalar için bir kontrol sisteminin tasarım ve geliştirmesine liderlik etti; proje maliyetlerini \textbf{\%30} azalttı ve \textbf{\$80k} tasarruf sağladı.}
      \end{cvitems}
    }

%---------------------------------------------------------
  \cventry
    {YAZILIM GELİŞTİRME MÜHENDİSİ} % Job title
    {} % Organization
    {} % Location
    {Tem 2019 - Ara 2021} % Date(s)
    {
      \begin{cvitems} % Description(s) of tasks/responsibilities
        \item {\textbf{IIoT} sistemlerini \textbf{Focas} kütüphanesi ile entegre ederek gelişmiş gerçek zamanlı görünürlükle kanıtlanan CNC izleme doğruluğunu \textbf{\%48} artırdı.}
        \item {\textbf{CNC}, \textbf{PLC} ve web servislerinden gelen sinyalleri analiz ederek \textbf{OEE} analiz araçlarını optimize etti, verimliliği \textbf{\%12} artırdı ve makine boşta kalma sürelerini minimize etti.}
      \end{cvitems}
    }

%---------------------------------------------------------
  \cventry
    {YARI ZAMANLI YAZILIM GELİŞTİRME MÜHENDİSİ} % Job title
    {} % Organization
    {} % Location
    {Tem 2018 - Tem 2019} % Date(s)
    {
      \begin{cvitems} % Description(s) of tasks/responsibilities
        \item {\textbf{PLC} ve \textbf{HMI} ile kaynak sürecini otomatikleştirerek, kısalan çevrim süresiyle gerçekleşen \textbf{3 kat daha hızlı üretim çevrimi} sağladı.}
      \end{cvitems}
    }

%---------------------------------------------------------
\end{cventries}

\begin{cventries}

%---------------------------------------------------------
  \cventry
    {YARI ZAMANLI PROJE MÜHENDİSİ} % Job title
    {DIMES} % Organization
    {İzmir, Türkiye} % Location
    {Tem 2017 - Haz 2018 } % Date(s)
    {
      \begin{cvitems} % Description(s) of tasks/responsibilities
        \item {\textbf{PLC} ve \textbf{IIoT} ile gerçek zamanlı izleme yaparak, bakım ihtiyaçlarındaki azalmayla doğrulanan verimliliği \textbf{\%4} artırdı.}
      \end{cvitems}
    }
    \vspace{1mm}

%---------------------------------------------------------
\end{cventries}
